\begin{established}[One order and precision for the two orientations]
\label{hyp:common-slice-scales}\label{par:slice-scales}
\hypertarget{source-slice-scales}{}
Take the selected rank of Lemma~\ref{lem:rank-window}. Write
\(L=x^\epsilon,\ b=e^L,\ z_0=frL,\ D_*=f\binom r2\).
With the exponents \(A_2=12,\ C_p=1\) of
Lemma~\ref{lem:window-parameter-choice}, put \(P=C_2z_0^{A_2}\)
for the fixed precision constant \(C_2\). On both runs, at every lift
and all sufficiently large selected scales, the order \(T=\lceil
x^{377}\rceil\) satisfies
\[
 \frac{P^4D_*^3kf\,x^{C_p}}{m\delta_2}\le T,\qquad
 P^3T^2\le\sqrt b,\qquad
 (D_*+2)n_W+2r\le T,\qquad 40(D_*+2),\,11D_*^2\le T .
\]
The cut has dimension parameter \(s=D_*\), and
\(n_P=8fr+4\le n_W\).
\label{cor:slice-scales-1}
\end{established}
\begin{proof}
For \(p+q=r\) the unitary base has dimension \(fpq\); the flags in
the two summands contribute \(f\binom p2+f\binom q2\).
Their sum is \(D_*\), including at \(q=0\).
Lemma~\ref{lem:rank-window} chooses
\[
 n_W=8\left\lceil\frac{10r}8
         \left(f+\frac{f^2(k-1)^3L^2}{m}\right)\right\rceil,
 \qquad r+t=n_P+Tn_W .
\]
In particular \(n_W\ge10fr\ge n_P\). No order is chosen anew after
the rank \(t\) has been fixed.

For an explicit comparison put \(y=x^{1/100}\). The bounds in that
lemma give \(f\ll y^{216}\), \(k,r\ll y^{470}\), \(m\ge y^{49}\),
\(L\le y^2\), and \(n_W\ll y^{2267}\). The nonempty set of larger
indices has density \(\delta_2\ge 1/n\), where \(2n=f\). Thus
\(z_0\ll y^{688}\), \(D_*\ll y^{1156}\), and
\[
 \frac{P^4D_*^3kf\,x}{m\delta_2}
 \ll_{C_2} y^{48\cdot688+3\cdot1156+470+216+100-49+216}
 =y^{37445}\ll x^{377}.
\]
The remaining lower bounds for \(T\) have still smaller exponents.
Finally \(\log(P^3T^2)=O_{C_2}(\log x)\), whereas
\(\log\sqrt b=x^\epsilon/2\). All quantities in this comparison are
fixed before the orientation and lift are chosen.
\end{proof}

\begin{lemma}[Size bounds in both orientations]
\label{lem:common-slice-scales}
For the data of Lemma~\ref{hyp:common-slice-scales}, both orientations satisfy the finite boundary estimate, the allowance conditions in \eqref{eq:slice-remainder-scales}, and the inclusion at every weight coordinate
\[
 T(w_L)_{\tau,i}\le\kappa_{\tau,i}\qquad\text{for all }\tau,i.
\]

For fixed choices of the constants in the polynomial \(P\), the
threshold on \(x\) is common to all lifts in the two runs.
\end{lemma}

\begin{proof}
\proofstep{Compatibility of the weight inclusions.}
The weight inclusion is coordinatewise:
\(T\) times the light rows is bounded by the selected
\(\kappa\) rows, so \(TD_L\le D_\kappa\).
The weight size has the differential degree, a central term
\(fr\mu_j n\), and a constant one, with \(1\le\mu_j\le r\).
The three central parameters at \(w_L,w_F,\kappa\) are
\(n_W,n_P,Tn_W+n_P\). The constant term in the weight size gives \(TS_L\le S_\kappa+T\).
Using \(T\le b\) and the head-weight bounds \(b\le S_\kappa\le7kfb\),
\[
 TS_L\le S_\kappa+T\le S_\kappa+b\le2S_\kappa .
\]

For the short weight, the row-gap
remainder after subtracting \(T w_L\) is at most
\((s+2)\) times the light gap plus \(3T\). For the
terminal gap of the large block the adjustment
\((s+2)n_W+2r\) is at most \(T\), giving that same remainder
bound. Summing the gap comparisons from the last row in
each block bounds the differential degree by
\((s+2)D_L+3fr^2T\). The short central term is at most
\(fr^2 n_P\le fr^2T\). Including the constant term gives
\(|w_F|\le(s+2)S_L+5fr^2T\).
\label{cor:slice-scales-2}

\proofstep{Room for the finite boundary term.}
For the finite boundary cost write \(Q=\min(g_0,g_1)\),
\(R_7=(\operatorname{rd}E)^7\). The record uses
\[
             c=R_7P D_L^2/Q,
             \qquad g_0^*=1250r .
\]
The level-size input is \(Q\ge(kb)^8\).
To see the large margin in this comparison, let
\(B=\sqrt b\). The rank and weight data give
\(r\le2k^2\), \(f,m\le k\) and \(D_L\le2frb\le4k^2B^2\),
where \(f\le k\) and \(r\le2k\) give the second bound.
Also \(R_7\le B\), \(P^3T^2\le B\), \(TP\le B\)
and \(B\ge10^{100}\). Thus
\[
\begin{split}
 R_7^4P^3D_L^2 &\le16 k^4 B^9,\\
 Tg_0^*R_7PD_L^2 &\le40000 k^6 B^6,\\
 2fm^3 R_7^3 P D_L^2 &\le32 k^8 B^8 .
\end{split}
\]
Each is at most \(k^8 B^{16}\le Q\).

The middle bound
gives \(Tc g_0^*\le1\), and the last pays the auxiliary
cusp allowance at the same order. The short-basis
comparison for the \hyperref[par:slice-finite-depth]{finite readings} uses
\(R_0\le R_7\) and \(N_0^4n_W+2N_0^3\le P\).
These give its two pointwise inequalities for \(c\).
Also \[
 J_0\le2\sqrt{N_0}\le3k,\qquad D_LJ_0\le12k^3 b<Q.
\] The degree estimate
\(2D_\kappa+2<g_i\), \(TD_L\le D_\kappa\) and
\(T\ge40(D_*+2)\) give the level conditions for the
Hodge--Tate comparison. The inequality \(4D_L+4\le g_i\)
also follows directly from \(D_L\le4k^2b\) and
\(g_i\ge(kb)^8\).
\label{cor:slice-scales-3}

\proofstep{Checking the four coefficient allowances.}
The coefficient bounds for the four prices use the
polynomial domination at this same \(z_0\), including
the auxiliary, height and quotient constants at their
fixed exponents. Its threshold is chosen once after
those constants have been chosen. The record supplies
\(z_0\ge x\) from \(x\le fr\) and \(L\ge1\), so a
single lower bound on \(x\) ensures the required
polynomial domination for every place and for all
lifts on either run.
In particular the single threshold \(x_*\) chosen for the comparison has the implication
\[
  x\ge x_*\quad\Longrightarrow\quad z_0\ge x_*.
\]
 We use the construction in the specified data on the common family, including the
level interval and positive density of the indices with larger weights, as a
size input here. This argument derives the two-run
inequalities used by the index on that record.
\label{cor:slice-scales-4}
\end{proof}
